% Structural facts for two actual increasing arms, not the strict-core theorem.
\section{Rising differences of the reversed ordinary two-arm product}
\begin{proposition}
Let $X,Y$ be increasing arms and
$p(z)=z^{a+b}F_X(1/z)F_Y(1/z)$, with $a=\operatorname{wt}X$
and $b=\operatorname{wt}Y$. Then $p$ is unimodal and log-concave
on its rising side. In addition, the positive part of its
first-difference sequence is unimodal.
\end{proposition}
\begin{proof}
Ordinary two-arm product unimodality is an established input;
the proof by interval layers and final-part deletion is reproduced
in the supplied foundation. Its strengthened midpoint assertion
says that, when $a=b$ and both factors are nonrectangular,
$(F_XF_Y)_{a+1}\ge(F_XF_Y)_a$.

Induct on the total number of arm parts for the two additional
properties. If either arm polynomial is an interval, apply the
uniform-window results to the reverse of the other arm. Otherwise
orient $a\ge b$ and write $X=Zt$. Reversal gives
\[
 p=[t]q_Zq_Y+z^a q_Y.
\]
Let $m_0$ be the first mode of $p$. First-half increase of the
ordinary product and its unimodality imply that its last mode
is at least $\lceil(a+b)/2\rceil$. Thus
$m_0\le\lfloor(a+b)/2\rfloor$. If $a>b$, this gives $m_0<a$.
If $a=b$, the strengthened midpoint assertion implies that the
last forward mode is at least $a+1$, and hence $m_0\le a-1$.
This reasoning explicitly allows a maximal plateau.

In both weight cases $m_0<a$. Up to and including
$m_0$, $p$ agrees with the uniform window $[t]q_Zq_Y$; the
added term begins at $a$. The positive first differences of
$p$ are therefore a prefix of those of this window, followed
by zeros. The window has unimodal positive first differences
by induction and the uniform-window lemma. This proves that
the positive first differences of $p$ are unimodal.

For log-concavity before the first mode, the same agreement
and the rising-side window lemma apply. A window cannot have
passed its first mode and then strictly increase to the first
mode of $p$; consequently its relevant prefix lies on its
rising side. Log-concavity at a mode is automatic, since both
neighbors are at most the mode value. Interior points of a
maximal plateau also satisfy equality. This proves rising-side
log-concavity of $p$, including its last mode.
\end{proof}

\paragraph{Priority and limitation.}
The rising-side log-concavity assertion is the reversal of the
ordinary-product tail log-concavity already reported in Chen's
public abstract; no priority is claimed for it. The proof above
is given to make its use independent of the inaccessible full
text. The positive first differences themselves need not be
log-concave: the exact counterexample recorded with the
interval-product difference theorem still applies. Neither
additional property alone proves the single crossing needed
after inserting the central peak.

\paragraph{Two stronger smoothing statements fail.}
Even an increasing arm need not be log-concave at all ranks
starting with rank two. For $Y=(1,2,2,2)$ the coefficients at
ranks $2,3,4$ are $2,3,5$, and $3^2<2\cdot5$.

Nor does a window as long as the strict peak make the positive
rising differences of the reversed ordinary product log-concave.
Take $X=1^{58}2^3$, $Y=1^{37}2^{13}$, and $M=3$.
The first mode of $p$ above is $60$. For the nonnegative sequence
$d_r=p_r-p_{r-1}$ on $0\le r\le60$, extended by zero,
the coefficients of $[3]d$ at ranks $55,56,57$ are
\[
 3194,\quad1974,\quad1220,
 \qquad1974^2-3194\cdot1220=-4.
\]
Thus even this proposed limited smoothing closure cannot be
used in the remaining coefficient bridge.
